\documentclass[11pt]{article}
\usepackage[utf8]{inputenc}
\usepackage{amsmath,amssymb}
\usepackage{hyperref}
\title{Robust Functional Genomics Screen under Distribution Shift}
\author{Anonymous}
\date{}
\begin{document}
\maketitle

\begin{abstract}
This paper studies Functional Genomics Screen. Grounded in a real, citable literature \cite{ref1,ref2,ref3,ref4,ref5,ref6,ref7,ref8},
we propose a focused method and report \textbf{illustrative} experiments.
All references below are real and verifiable (OpenAlex/arXiv); the reported
numbers are simulated to illustrate intended behaviour.
\end{abstract}

\section{Introduction}
Functional Genomics Screen is an active research area. This work targets a concrete gap identified from
the recent literature and contributes a method together with an evaluation protocol.

\section{Related Work (real literature)}
The following works, retrieved from public bibliographic APIs, frame our study:
\begin{itemize}
\item Carpenter et al. (2006) — "CellProfiler: image analysis software for identifying and quantifying cell phenotypes", Genome biology (cited 5520).
\item Hasegawa et al. (2000) — "P<scp>LANT</scp>C<scp>ELLULAR AND</scp>M<scp>OLECULAR</scp>R<scp>ESPONSES TO</scp>H<scp>IGH</scp>S<scp>ALINITY</scp>", Annual Review of Plant Physiology and Plant Molecular Biology (cited 4594).
\item Giaever et al. (2002) — "Functional profiling of the Saccharomyces cerevisiae genome", Nature (cited 4537).
\item Redon et al. (2006) — "Global variation in copy number in the human genome", Nature (cited 4383).
\item Winkel (2001) — "Flavonoid Biosynthesis. A Colorful Model for Genetics, Biochemistry, Cell Biology, and Biotechnology", PLANT PHYSIOLOGY (cited 3735).
\item Konermann et al. (2014) — "Genome-scale transcriptional activation by an engineered CRISPR-Cas9 complex", Nature (cited 2967).
\end{itemize}

\section{Method}
We model distribution shift explicitly and optimise a worst-case objective over an uncertainty set of plausible test distributions. We instantiate this idea for Functional Genomics Screen and analyse its cost/quality trade-off.

\section{Experiments (Illustrative)}
\begin{center}
\begin{tabular}{lcc}
System & Primary Metric & Efficiency \\ \hline
Strong baseline & 0.812 & 1.00$\times$ \\
Ablation & 0.828 & 0.92$\times$ \\
\textbf{Ours} & \textbf{0.851} & \textbf{0.71$\times$}
\end{tabular}
\end{center}
\emph{Metrics simulated to illustrate the intended effect; not measured.}

\section{Conclusion}
We connected a real literature to a concrete method for Functional Genomics Screen. Future work will
validate the illustrative results with real experiments.

\begin{thebibliography}{9}
\bibitem{ref1} Anne E. Carpenter, Thouis R. Jones, Michael R. Lamprecht et al.. CellProfiler: image analysis software for identifying and quantifying cell phenotypes. \emph{Genome biology}, 2006. DOI:10.1186/gb-2006-7-10-r100
\bibitem{ref2} Paul M. Hasegawa, Ray A. Bressan, Jian‐Kang Zhu et al.. P<scp>LANT</scp>C<scp>ELLULAR AND</scp>M<scp>OLECULAR</scp>R<scp>ESPONSES TO</scp>H<scp>IGH</scp>S<scp>ALINITY</scp>. \emph{Annual Review of Plant Physiology and Plant Molecular Biology}, 2000. DOI:10.1146/annurev.arplant.51.1.463
\bibitem{ref3} Guri Giaever, Angela Chu, Li Ni et al.. Functional profiling of the Saccharomyces cerevisiae genome. \emph{Nature}, 2002. DOI:10.1038/nature00935
\bibitem{ref4} Richard Redon, Shumpei Ishikawa, Karen Fitch et al.. Global variation in copy number in the human genome. \emph{Nature}, 2006. DOI:10.1038/nature05329
\bibitem{ref5} A. Winkel. Flavonoid Biosynthesis. A Colorful Model for Genetics, Biochemistry, Cell Biology, and Biotechnology. \emph{PLANT PHYSIOLOGY}, 2001. DOI:10.1104/pp.126.2.485
\bibitem{ref6} Silvana Konermann, Mark D. Brigham, Alexandro E. Trevino et al.. Genome-scale transcriptional activation by an engineered CRISPR-Cas9 complex. \emph{Nature}, 2014. DOI:10.1038/nature14136
\bibitem{ref7} Mathew J. Garnett, Elena J. Edelman, Sonja J. Heidorn et al.. Systematic identification of genomic markers of drug sensitivity in cancer cells. \emph{Nature}, 2012. DOI:10.1038/nature11005
\bibitem{ref8} Alan L. Harvey, RuAngelie Edrada‐Ebel, Ronald J. Quinn. The re-emergence of natural products for drug discovery in the genomics era. \emph{Nature Reviews Drug Discovery}, 2015. DOI:10.1038/nrd4510
\end{thebibliography}
\end{document}
